\section{Resilience in Management}

\subsection{Resilience as a Dynamic Managerial Capability}

Among the many domains contributing to resilience thinking, management plays a pivotal role in operationalizing and institutionalizing resilience across complex organizational landscapes. Within this domain, resilience is not merely a theoretical construct but a lived organizational reality that is shaped through practices, routines, and governance arrangements. Management research has gradually shifted from conceptualizing resilience as a reactive outcome to treating it as a proactive, path-dependent capability. This evolution reflects the broader theoretical progression in the field—from classical contingency-based models to more contemporary views rooted in dynamic capabilities, complex systems, and institutional embeddedness.\\ 

Organizational resilience is increasingly framed as a dynamic capability that allows firms to sense, anticipate, respond to, and even leverage disruptions for transformation \autocite{duchek2020organizational, teece2007explicating}. The dynamic capabilities framework emphasizes the processes by which firms purposefully create, extend, or modify their resource base to meet environmental challenges. Within this framing, resilience emerges from microfoundational elements, including routines, decision-making heuristics, and cognitive schemas that shape how actors interpret and respond to uncertainty \autocite{hepfer2022heterogeneity}. It is not simply about robustness or risk mitigation; rather, it is an adaptive orientation embedded in the fabric of strategic and operational activity. The interplay between slack resources, distributed leadership, and organizational identity often underpins this capacity, enabling firms to deploy resources flexibly and sustain performance in the face of adversity \autocite{bhamra2011resilience, barasa2018what}. In practice, this can manifest through investment in absorptive capacity, dual structures that promote both innovation and efficiency, and intentional resilience audits that assess an organization's readiness to absorb and adapt to disruptive shocks.

\subsection{Design, Culture, and the Construction of Resilience}
From an organizational design perspective, resilience is fostered through both structural and cultural mechanisms. Structural enablers include modular architectures, redundancy in critical functions, and decentralization of authority to support rapid decision-making \autocite{lengnick2011developing}. These allow systems to reconfigure rapidly and isolate failures, limiting their propagation across units. Culturally, resilience is nurtured through norms that value learning, reflexivity, and psychological safety. A resilient culture embraces ambiguity, empowers employees to experiment, and promotes distributed sensemaking across hierarchical levels. Leadership plays a critical role in this domain, as resilient organizations often display leadership that is forward-looking, emotionally intelligent, and capable of mobilizing cross-functional alignment \autocite{gillard2023efficient}.\\

Participatory mechanisms—such as hackathons, cross-disciplinary task forces, or innovation labs—further reinforce this adaptive culture by facilitating the emergence of novel solutions and strengthening collective sensemaking \autocite{maillart2024computational}. These configurations also illustrate how resilience is co-constructed through iterative negotiation among actors, practices, and technologies. Over time, such practices may crystallize into organizational routines that serve as repositories of institutional memory and adaptive learning. Together, these structural and cultural elements underscore that resilience is not exogenous or incidental but instead deeply embedded in how an organization configures its resources, cultivates its people, and governs its practices.

\subsection{Cybersecurity Resilience as a Strategic and Institutional Imperative}

The cybersecurity domain serves as a particularly illustrative context in which these resilience dynamics unfold. Cyber risks are characterized by high uncertainty, rapid evolution, and interdependent vulnerabilities that cut across technical, human, and organizational layers \autocite{linkov2013resilience}. Unlike traditional risk scenarios that are probabilistic and often bounded in scope, cybersecurity threats are emergent, asymmetrical, and often invisible until activated. Traditional risk management approaches, focused on compliance and control, are often inadequate in the face of such complexity. Consequently, cyber resilience is now viewed as an integrated capability encompassing not only technological safeguards but also strategic foresight, organizational agility, and institutional alignment \autocite{bjorck2015cyber, lostri2022shared}.\\

Frameworks such as ISO 22301 provide structured guidance for crisis response, emphasizing the importance of maintaining core functions under stress \autocite{iso201922301}. However, emergent views stress the need for adaptive processes rooted in learning and reflexive governance \autocite{hausken2020cyber}, where security practices evolve in tandem with threat landscapes. Regulatory pressures—such as those stemming from Basel III—have institutionalized resilience as a strategic imperative, linking it to organizational legitimacy and systemic risk management \autocite{baselcommittee2011basel, almeida2020challenges}. This has catalyzed a broader governance turn in cybersecurity management, where firms are increasingly accountable not only for their individual risk posture but also for their contribution to systemic resilience.\\

Meanwhile, advancements in AI, predictive analytics, and scenario modeling empower firms to anticipate cyber threats and orchestrate adaptive responses in real time \autocite{kuypers2018designing, gillard2023efficient}. These technologies allow for the simulation of threat vectors, visualization of attack surfaces, and real-time adjustment of defensive protocols. In sum, cyber resilience exemplifies the synthesis of technological infrastructure, human cognition, and institutional design --demonstrating how resilience in management is both a means of navigating uncertainty and a strategy for sustained adaptation. When viewed through the lens of resilience, cybersecurity becomes less about perimeter defense and more about the dynamic coordination of resources, people, and information to enable anticipatory and responsive action at scale.